\scanpage{72}
\mspage{G.}
% archive foliation: 71
diviseur $\sum n_iX_i$, où les $X_i$ sont les ``diviseurs singuliers''
pour $f$. Il reste donc à voir que si une telle combinaison linéaire est
\deleted{algébriquement}
$\tau$-équivalente à $0$, elle est nulle.
\deleted{\unsure{Or}}
C'est là un lemme bien connu \ldots

\margin{vérifier ? On peut se ramener au cas projectif}
\margin{\uline{insuffisant}}

Ainsi, $f:\underline{\mathrm{Pic}}_Y^\tau\longrightarrow
\underline{\mathrm{Pic}}_X^\tau$ est une \uline{immersion}
\uline{surjective}.
Il \uline{reste à prouver} qu'elle induit un isom sur les espaces tangents
à l'origine, donc à vérifier le

\uline{Corollaire} Sous les conditions du th., l'homomorphisme
\[
f^*:H^1(Y,\underline O_Y)\longrightarrow H^1(X,\underline O_X)
\]
est un isomorphisme.

On sait déjà que c'est injectif, la question est la surjectivité.
\deleted{En effet}
En caractéristique $p>0$, il suffit de prouver qu'on a bijectivité
\[
H^1(Y,\underline O_Y)^F
\mathbin{\unsure{-|\longrightarrow}}
H^1(X,\underline O_X)^F
\]
\[
{}_{\unsure{F}}H^1(Y,\underline O_Y)
\mathbin{\unsure{-|\longrightarrow}}
{}_{\unsure{F}}H^1(X,\underline O_X)
\]
où $F$ est le Frobenius.
Le premier hom s'écrit
\[
H^1(Y,\mathbf Z/p\mathbf Z)\longrightarrow H^1(X,\mathbf Z/p\mathbf Z)
\]
